% Generated draft: requires mathematical and source review.
\documentclass[11pt]{article}
\usepackage[margin=25mm]{geometry}
\usepackage{fontspec}
\setmainfont{Latin Modern Roman}
\usepackage{amsmath,amssymb,mathtools,unicode-math}
\setmathfont{Latin Modern Math}
\usepackage{enumitem,needspace,xurl,hyperref}
\setlength{\parindent}{0pt}
\setlength{\parskip}{5pt}
\setlength{\emergencystretch}{3em}
\newcommand{\N}{\mathbb{N}}
\newcommand{\Z}{\mathbb{Z}}
\newcommand{\Q}{\mathbb{Q}}
\newcommand{\R}{\mathbb{R}}
\newcommand{\C}{\mathbb{C}}
\newcommand{\F}{\mathbb{F}}
\newcommand{\A}{\mathbb{A}}
\newcommand{\PP}{\mathbb{P}}
\newcommand{\E}{\mathbb{E}}
\newcommand{\eps}{\varepsilon}
\newcommand{\dd}{\,\mathrm{d}}
\begin{document}
\section*{UnsolvedMath ID: 800404}
\textbf{Source problem code:} catalogue entry
\subsection{Definitions and mathematical statement}
Let \(f\colon\PP^N\dashrightarrow\PP^N\) be a dominant rational map, and let \(f^n\) denote its \(n\)-fold iterate for each positive integer \(n\). Here dominance means that the image is Zariski dense. The ground field and its characteristic are not specified in the supplied record; the formulation below requires the intended ground field to be fixed.

Write \(\deg(f^n)\) for the algebraic degree of the iterate, computed from homogeneous coordinate polynomials of a common degree after removing common factors. Interpret the otherwise undefined notation \(\delta_f\) as the first dynamical degree,
\[
\delta_f=\lim_{n\to\infty}\deg(f^n)^{1/n}.
\]
This interpretation of \(\delta_f\) requires confirmation against the source context.

For positive sequences \(F,G\), the notation \(F(n)\asymp G(n)\) means that there exist positive constants \(c_1,c_2,c_3\) such that
\[
c_1F(n)\le G(n)\le c_2F(n)
\quad\text{for every integer }n\ge c_3.
\]

The supplied question asks whether, for every such \(f\), there exists an integer \(\ell_f\) with \(0\le\ell_f\le N\) such that
\[
\deg(f^n)\asymp n^{\ell_f}\delta_f^n.
\]
Explicitly, does there exist such an \(\ell_f\), together with positive constants \(c_1,c_2,c_3\), all allowed to depend on \(f\) but not on \(n\), for which
\[
c_1\deg(f^n)\le n^{\ell_f}\delta_f^n\le c_2\deg(f^n)
\quad\text{for every integer }n\ge c_3?
\]
The supplied registry labels this question open; that label is not an independently verified assessment of its current mathematical status.
\subsection{Short English statement}
Does the degree of every iterate of a dominant rational self-map of \(\PP^N\) grow, up to fixed positive multiplicative constants, like \(n^{\ell_f}\delta_f^n\) for some integer \(0\le\ell_f\le N\)? The ground field is unspecified, and interpreting \(\delta_f\) as the first dynamical degree requires source-context confirmation.
\subsection{Sources}
\begin{itemize}
\item Local UnsolvedMath registry, permanent ID \(800404\), exact source problem code \texttt{catalogue entry}, title ``Open Problems in Arithmetic Dynamics --- Question 27.'' The supplied record labels the problem open.
\item Benedetto, DeMarco, Ingram, Jones et al., \emph{Current Trends and Open Problems in Arithmetic Dynamics} (2018), Question 27. Supplied source URL: \url{https://arxiv.org/abs/1806.04980}. The registry attributes the problem to these authors and gives the proposed year as 2018.
\item Supplied provenance reports access on 2026-10-16 and extraction by \texttt{arxiv\_latex}. These are supplied provenance claims; neither access to the source nor its current status has been independently verified here.
\end{itemize}
\end{document}
